\documentclass[12pt,a4paper]{article}
\usepackage[utf8]{inputenc}
\usepackage{amsmath}
\usepackage{amssymb}
\usepackage{amsfonts}
\usepackage{graphicx}
\usepackage{hyperref}
\usepackage{geometry}
\geometry{margin=2.5cm}

\title{BEC-Objekte: Die kompaktesten Quellen der Materie\\ 
\large{Eine konsistente Beschreibung aus der WDBT+}}
\author{Michael Czybor}
\date{\today}

\begin{document}

\maketitle

\begin{abstract}
Die Weber-de-Broglie-Bohm-Theorie mit fraktalem Raum (WDBT+) liefert ein konsistentes, singularitätenfreies Modell der Gravitation. In diesem Rahmen werden kompakte Objekte – in der Allgemeinen Relativitätstheorie als Schwarze Löcher bezeichnet – nicht als Materiesenken, sondern als Quellen kontinuierlicher Materieentstehung interpretiert. Die Dynamische Schwere-Trägheits-Theorie (DSTT) erzwingt mit $\dot{\beta}>0$ eine universelle Auswärtsdrift, die zur Bildung von Galaxien führt.

Im Zentrum dieser Beschreibung stehen \textbf{Bose-Einstein-Kondensat (BEC)-Objekte}: die kompaktesten Objekte im Universum. Sie entstehen durch Paarbildung von Fermionen und werden durch das Quantenpotential stabilisiert. Ihre minimale Größe ist durch die fundamentale Längenskala $l_0 \approx 1{,}8 \times 10^{-15}\,\text{m}$ begrenzt, die aus der fraktalen Raumdimension $D \approx 2{,}71$ folgt. Daraus ergibt sich eine \textbf{maximale Masse} für ein einzelnes BEC-Objekt von 
\[
M_{\text{BEC}} \approx 3 \times 10^6 M_\odot.
\]

Beobachtete supermassereiche Objekte (bis $10^{10} M_\odot$) sind keine einzelnen kompakten Objekte, sondern \textbf{Systeme} aus einem zentralen BEC-Objekt und einer massiven Umgebung (Akkretionsscheibe, Sternhaufen, Gas). Diese Umgebung entsteht durch die kontinuierliche Materieentstehung im Quantenpotential des Kerns und die anschließende Drift nach außen. Das Verhältnis von Galaxienmasse zu Kernmasse ($M_{\text{Gal}}/M_{\text{Kern}} \approx 1000$) folgt aus der fraktalen Raumdimension.

Bei Kollisionen verhalten sich BEC-Objekte wie normale Objekte: Sie können verschmelzen (wenn die Gesamtmasse $M_{\text{BEC}}$ nicht überschreitet), abprallen oder zerreißen. Dies führt zu Doppelkern-Systemen oder zur Überführung von Masse in die Umgebung.

Das Modell erklärt den Sonnenwind ohne signifikanten Masseverlust der Sonne, die flachen Rotationskurven von Galaxien ohne Dunkle Materie sowie die natürliche Obergrenze für kompakte Objekte. Es kommt ohne Urknall, Dunkle Energie und Singularitäten aus und basiert ausschließlich auf den Grundlagen der Weber-Gravitation, des Quantenpotentials und der fraktalen Raumstruktur.
\end{abstract}

\section{Einleitung}

Die Allgemeine Relativitätstheorie (ART) beschreibt Gravitation als Krümmung der Raumzeit. Sie hat sich in vielen Experimenten bewährt, führt jedoch zu konzeptionellen Problemen: Singularitäten im Zentrum Schwarzer Löcher und am Urknall, die Notwendigkeit von Dunkler Materie und Dunkler Energie sowie das Informationsverlustparadoxon. Die WDBT+ bietet eine alternative, fundamentalere Beschreibung, aus der die ART als emergenter Grenzfall hervorgeht.

\subsection{Grundlagen der WDBT+}

Die WDBT+ basiert auf drei Säulen:

\begin{enumerate}
\item \textbf{Weber-Gravitation (WG)}: Eine geschwindigkeits- und beschleunigungsabhängige Kraft
\[
\vec{F}_{\text{WG}} = -\frac{GMm}{r^2} \left\{ \left[ 1 - \frac{\dot{r}^2}{c^2} + \beta_{\text{WG}} \frac{r\ddot{r}}{c^2} \right] \hat{r} + \frac{2(\hat{r} \cdot \vec{v})}{c^2} \vec{v} \right\}
\]
mit $\beta_{\text{WG}}=0{,}5$ für Massen.

\item \textbf{de Broglie-Bohm-Quantenpotential (DBT)}:
\[
Q = -\frac{\hbar^2}{2m} \frac{\nabla^2 R}{R}
\]
Dieses Potential verhindert Singularitäten und ermöglicht Materieentstehung aus dem Vakuum.

\item \textbf{Fraktale Raumdimension}:
\[
D = \frac{\ln 20}{\ln(2 + \phi)} \approx 2{,}71, \quad \phi = \frac{1+\sqrt{5}}{2}
\]
Sie liefert natürliche Cutoffs und erklärt die Feinstrukturkonstante.
\end{enumerate}

\subsection{Dynamische Schwere-Trägheits-Theorie (DSTT)}

Die DSTT trennt Ursache (Schwere) und Wirkung (Trägheit). Die Zustandsvariable
\[
\beta_{\text{DSTT}}(t) = \frac{\frac{1}{2}m r^2 \dot{\phi}^2}{\frac{1}{2}m \dot{r}^2 + \frac{1}{2}m r^2 \dot{\phi}^2}
\]
erfüllt $\dot{\beta}_{\text{DSTT}} > 0$ und erzwingt damit logarithmische Spiralbahnen
\[
r(\phi) = Ke^{\phi/B}
\]
und eine universelle Auswärtsdrift.

\section{BEC-Objekte: Die kompaktesten Objekte im Universum}

\subsection{Die fundamentale Längenskala $l_0$}

Aus der fraktalen Raumstruktur mit $D \approx 2{,}71$ folgt eine fundamentale Längenskala:
\[
l_0 = \left( \frac{V_{\text{Dodekaeder}}}{V_{\text{Einheitskugel}}} \right)^{1/3} \lambda_p \approx 1{,}8 \times 10^{-15}\,\text{m},
\]
wobei $\lambda_p$ die Compton-Wellenlänge des Protons ist. Diese Länge ist die \textbf{kleinste physikalische Distanz} – nichts kann kleiner sein. Sie ersetzt die Planck-Skala als natürliche untere Grenze.

\subsection{Paarbildung und Bose-Einstein-Kondensation}

Bei hinreichend hohen Dichten bilden Fermionen Cooper-Paare und werden zu effektiven Bosonen. Diese können in einen \textbf{Bose-Einstein-Kondensat (BEC)-Zustand} übergehen. In diesem Zustand:
\begin{itemize}
\item Entfällt der Fermi-Druck (der sonst weitere Kompression verhindert)
\item Das Quantenpotential dominiert die repulsive Wechselwirkung
\item Das Objekt kann bis zur fundamentalen Länge $l_0$ schrumpfen
\end{itemize}

\subsection{Größe und maximale Masse eines BEC-Objekts}

Im Gleichgewicht zwischen Gravitation und Quantenpotential ergibt sich für ein BEC-Objekt:
\[
R(M) = K \cdot M^{-0{,}91}
\]
Diese Funktion fällt mit steigender Masse. Sie erreicht die minimale Größe $R = l_0$ bei:
\[
M_{\text{BEC}} = \left( \frac{K}{l_0} \right)^{1/0{,}91} \approx 3 \times 10^6 M_\odot
\]

\textbf{Dies ist die maximale Masse eines einzelnen BEC-Objekts.} Oberhalb dieser Grenze wäre $R < l_0$, was physikalisch unmöglich ist.

\subsection{Was sind beobachtete supermassereiche Objekte?}

Beobachtet werden zentrale Objekte in Galaxien mit Massen bis $10^{10} M_\odot$. Diese sind \textbf{keine einzelnen kompakten Objekte}, sondern \textbf{Systeme} bestehend aus:
\begin{itemize}
\item Einem zentralen BEC-Objekt mit $M = M_{\text{BEC}} \approx 3 \times 10^6 M_\odot$ und $R = l_0$
\item Einer massiven Umgebung (Akkretionsscheibe, Sternhaufen, Gas), die die zusätzliche Masse von $10^9 - 10^{10} M_\odot$ enthält
\end{itemize}

Diese Umgebung entsteht durch die kontinuierliche Materieentstehung aus dem Quantenpotential des Kerns und die anschließende Drift nach außen.

\subsection{Strukturelle Analogie zum Atom}

Die beschriebene Struktur – ein winziger, extrem dichter Kern, umgeben von einer riesigen, dünnen Hülle, mit viel Vakuum dazwischen – erinnert an ein Atom. Im Atom sitzt fast die gesamte Masse im Kern ($\sim 10^{-15}\,\text{m}$), die Elektronen sind weit draußen ($\sim 10^{-10}\,\text{m}$). In der BEC-Galaxie ist der Kern ebenfalls winzig ($l_0 \approx 10^{-15}\,\text{m}$), die Umgebung (Sterne, Gas) erstreckt sich über $10^{21}\,\text{m}$ – und der Kern trägt nur einen kleinen Teil der Gesamtmasse.

Die Analogie ist nicht zufällig: In beiden Fällen organisiert ein kompaktes Zentrum eine ausgedehnte Umgebung. Im Atom geschieht dies durch elektromagnetische Wechselwirkung und Quantenmechanik, in der Galaxie durch Weber-Gravitation, Quantenpotential und die DSTT-Drift. Die fraktale Raumdimension $D \approx 2{,}71$ verbindet beide Welten – sie bestimmt sowohl die Feinstrukturkonstante als auch das Massenverhältnis $M_{\text{Gal}}/M_{\text{Kern}} \approx 1000$.

Dieser strukturelle Gleichklang deutet auf ein tieferes Prinzip hin: Die Natur organisiert sich wiederholt in hierarchischen Systemen mit einem aktiven Zentrum und einer ausgedehnten, durch das Zentrum erzeugten Peripherie.

\subsection{Die maximale Masse als natürliche Grenze}

Ein BEC-Objekt kann nicht über $M_{\text{BEC}}$ hinauswachsen. Wenn zusätzliche Materie durch Akkretion oder Verschmelzung hinzukommt:
\begin{itemize}
\item Die Überschussmasse kann \textbf{nicht ins BEC-Objekt integriert werden}
\item Sie verbleibt in der \textbf{Umgebung} (Akkretionsscheibe, Sternhaufen)
\item Das BEC-Objekt selbst bleibt bei $M_{\text{BEC}}$
\end{itemize}

Dies erklärt, warum die beobachteten Massen supermassereicher Objekte so stark streuen: Sie reflektieren die Masse der Umgebung, nicht die des Kerns.

\section{Kollisionen von BEC-Objekten}

\subsection{Verhalten wie normale Objekte}

BEC-Objekte haben keine Sonderstellung. Bei Kollisionen verhalten sie sich wie alle anderen physikalischen Objekte:

\begin{itemize}
\item \textbf{Verschmelzung}: Wenn $M_1 + M_2 \leq M_{\text{BEC}}$, können sie zu einem neuen BEC-Objekt verschmelzen.
\item \textbf{Abprallen}: Wenn die Relativgeschwindigkeit hoch genug ist, prallen sie elastisch auseinander.
\item \textbf{Zerrüttung}: Bei geringer Relativgeschwindigkeit und $M_1 + M_2 > M_{\text{BEC}}$ können die Kerne zerreißen; die Masse wird in die Umgebung überführt.
\item \textbf{Doppelkern-Systeme}: Die Kerne bleiben getrennt und umkreisen sich in einer gemeinsamen Umgebung.
\end{itemize}

\subsection{Gravitationswellen}

Verschmelzungen von BEC-Objekten erzeugen Gravitationswellen – ähnlich wie Neutronensterne oder Schwarze Löcher in der ART. Unterschiede:

\begin{itemize}
\item BEC-Objekte haben \textbf{keinen Ereignishorizont}, sondern eine Oberfläche (auch wenn diese extrem klein ist)
\item Die Wellenform kann abweichen, weil die Materie nicht "verschluckt" wird
\item Nach der Verschmelzung kann es \textbf{Nachglühen} geben, wenn Materie aus den Umgebungen akkretiert wird
\end{itemize}

Die LIGO/Virgo-Daten sind mit beiden Modellen konsistent, da die beobachteten Massen ($\sim 10-100 M_\odot$) weit unter $M_{\text{BEC}}$ liegen.

\section{Materieentstehung und Galaxienbildung}

\subsection{Kontinuierliche Materieentstehung}

Das Quantenpotential ist im Vakuum nicht null und kann Arbeit verrichten. Die Materieerzeugungsrate ist proportional zu $Q$:
\[
\dot{\rho}_{\text{cre}}(r) \propto Q(r) \propto \frac{1}{r^2}.
\]
Im fraktalen Raum mit $D \approx 2{,}71$ modifiziert sich die radiale Verteilung:
\[
\frac{d\dot{M}_{\text{cre}}}{d\ln r} \propto r^{3-D} = r^{0{,}29}.
\]

\subsection{Drift und Einfang}

Die DSTT liefert die Driftgeschwindigkeit
\[
v_{\text{drift}}(r) = B \sqrt{\frac{GM}{r^3}},
\]
die Fluchtgeschwindigkeit ist
\[
v_{\text{esc}}(r) = \sqrt{\frac{2GM}{r}}.
\]
Der kritische Radius $r_c$, bei dem beide gleich sind, folgt aus
\[
B \sqrt{\frac{GM}{r_c^3}} = \sqrt{\frac{2GM}{r_c}} \quad\Rightarrow\quad r_c = \frac{B}{\sqrt{2}}.
\]

\subsection{Einfangwahrscheinlichkeit}

Im fraktalen Raum ist die Einfangwahrscheinlichkeit
\[
\alpha_{\text{in}} = \frac{\int_{r_{\text{min}}}^{r_c} r^{3-D} \, d\ln r}{\int_{r_{\text{min}}}^{r_{\text{max}}} r^{3-D} \, d\ln r}
\approx \left( \frac{r_c}{r_{\text{max}}} \right)^{3-D},
\]
wobei $r_{\text{max}} \sim \sqrt[3]{GM/\rho_{\text{vac}}}$ der Einflussbereich ist ($\rho_{\text{vac}}$ ist die Vakuumenergiedichte). Mit $D \approx 2{,}71$ ist $3-D = 0{,}29$.

\subsection{Galaxienentstehung aus dem BEC-Kern}

Für ein BEC-Objekt mit $M = M_{\text{BEC}} \approx 3 \times 10^6 M_\odot$ ergibt sich:
\[
r_{\text{max}} \sim 5 \times 10^{11}\,\text{m},\qquad
r_c \sim 200\,\text{m},\qquad
\alpha_{\text{in}} \approx 0{,}001.
\]

Daraus folgt das Verhältnis von Galaxienmasse zu Kernmasse:
\[
\frac{M_{\text{Gal}}}{M_{\text{Kern}}} = \frac{1 - \alpha_{\text{in}}}{\alpha_{\text{in}}} \approx 1000.
\]

Dies stimmt mit Beobachtungen überein. Die Galaxie entsteht aus der ausströmenden Materie, die im Lauf der Zeit akkumuliert wird. Die logarithmische Spirale der DSTT erklärt die Spiralstruktur; die flachen Rotationskurven folgen aus der DSTT-Dynamik ohne Dunkle Materie.

\section{Die Sonne als Mikrokosmos}

Das Modell ist skalierbar. Für $M = M_\odot$:
\[
r_{\text{max}} \sim 1{,}5 \times 10^{13}\,\text{m},\qquad
r_c \sim 1{,}4 \times 10^{-11}\,\text{m},\qquad
\alpha_{\text{in}} \approx 1{,}2 \times 10^{-7}.
\]

\subsection{Interpretation des Sonnenwinds}

Im Standardmodell wird der Sonnenwind als Materieausstoß aus der Sonnenkorona gesehen. In der WDBT+ hingegen:
\begin{itemize}
\item Materie entsteht in der Nähe der Sonne (bei $r \sim 10^{-11}$ m)
\item Die DSTT-Drift treibt sie nach außen
\item Nur ein winziger Bruchteil ($\sim 10^{-7}$) wird von der Sonne eingefangen
\item Die Sonne verliert daher keine signifikante Masse
\end{itemize}

Der beobachtete Sonnenwind-Massestrom von $\sim 10^{-14} M_\odot/\text{Jahr}$ ist konsistent mit dieser Interpretation.

\subsection{Testbare Vorhersagen}

\begin{enumerate}
\item \textbf{Konstanz der Sonnenmasse}: Präzise Messungen (z.B. mit Raumsonde BepiColombo) sollten zeigen, dass die Sonne keinen messbaren Masseverlust durch den Sonnenwind erfährt.
\item \textbf{Isotopenanomalien}: Neu entstehende Materie könnte leichte Isotope bevorzugen; solche Anomalien im Sonnenwind wären ein Fingerabdruck der Materieentstehung.
\item \textbf{Fraktale Skalierung}: Die Fluktuationen des Sonnenwinds sollten ein Skalierungsverhalten mit Exponent $0{,}29$ zeigen, das im Standardmodell nicht vorkommt.
\item \textbf{Maximale Kernmasse}: Die beobachteten zentralen Objekte in Galaxien sollten eine natürliche Obergrenze bei $\sim 3 \times 10^6 M_\odot$ haben – darüber liegende Massen sind Umgebungsmassen.
\item \textbf{Doppelkern-Systeme}: Bei Galaxienverschmelzungen sollten Doppelkerne häufiger sein als in der ART vorhergesagt, da BEC-Objekte mit $M_1 + M_2 > M_{\text{BEC}}$ nicht verschmelzen.
\end{enumerate}

\section{Vergleich mit dem Standardmodell}

\begin{table}[h]
\centering
\begin{tabular}{|l|c|c|}
\hline
\textbf{Phänomen} & \textbf{Standardmodell} & \textbf{WDBT+} \\
\hline
Zentrale Objekte & Senken (verschlingen Materie) & Quellen (erzeugen Materie) \\
Materiefluss & Einwärts (Akkretion) & Auswärts (Drift) \\
Galaxien & Kollaps + Dunkle Materie & Wachstum von innen nach außen \\
Universum & Expandierend, Urknall & Stationär, ewig \\
Schwarze Löcher & Singularität + Horizont & BEC-Objekt, kein Horizont \\
Maximale Masse (Kern) & Keine natürliche Grenze & $M_{\text{BEC}} \approx 3 \times 10^6 M_\odot$ \\
Super-massereich & Einzelnes Schwarzes Loch & BEC-Kern + Umgebung \\
Kollisionen & Immer Verschmelzung & Wie normale Objekte \\
Sonnenwind & Ausstoß aus der Korona & Neu entstehende Materie \\
Dunkle Materie & Benötigt & Nicht benötigt \\
Dunkle Energie & Benötigt & Nicht benötigt \\
\hline
\end{tabular}
\caption{Gegenüberstellung der Standardkosmologie und der WDBT+.}
\end{table}

\section{Kosmologische Konsequenzen}

Die WDBT+ führt zu einem stationären, nicht-expandierenden Universum (Steady State). Die beobachtete Rotverschiebung ist keine Folge der Expansion, sondern kumulative Gravitation. Das Quantenpotential ermöglicht kontinuierliche Materieentstehung, die im stationären Universum die beobachtete Struktur erhält.

Die JWST-Daten, die im Standardmodell als "zu frühe Galaxien" erscheinen, sind im Rahmen der WDBT+ keine Überraschung: Sie zeigen alte Galaxien in einem ewigen Universum.

\section{Zusammenfassung}

Die WDBT+ liefert ein konsistentes, singularitätenfreies Bild der Gravitation:

\begin{itemize}
\item Die kompaktesten Objekte im Universum sind \textbf{BEC-Objekte} – Bose-Einstein-Kondensate aus Cooper-Paaren, stabilisiert durch das Quantenpotential.
\item Ihre minimale Größe ist die fundamentale Längenskala $l_0 \approx 1{,}8 \times 10^{-15}\,\text{m}$, ihre maximale Masse $M_{\text{BEC}} \approx 3 \times 10^6 M_\odot$.
\item Beobachtete supermassereiche Objekte (bis $10^{10} M_\odot$) sind \textbf{Systeme} aus einem zentralen BEC-Kern und einer massiven Umgebung (Akkretionsscheibe, Sternhaufen, Gas).
\item Bei Kollisionen verhalten sich BEC-Objekte wie normale Objekte: Sie können verschmelzen ($M_1+M_2 \leq M_{\text{BEC}}$), abprallen oder zerreißen.
\item Das Quantenpotential erzeugt kontinuierlich neue Materie, bevorzugt in der Nähe existierender Masse.
\item Die DSTT erzwingt eine universelle Auswärtsdrift ($\dot{\beta}>0$), die zur Bildung von Galaxien führt.
\item Das Verhältnis $M_{\text{Gal}}/M_{\text{Kern}} \approx 1000$ folgt aus der fraktalen Raumdimension $D\approx2{,}71$.
\item Das Modell ist skalierbar und erklärt den Sonnenwind ohne Masseverlust der Sonne.
\item Es kommt ohne Dunkle Materie, Dunkle Energie, Urknall und Singularitäten aus.
\end{itemize}

Die Theorie macht testbare Vorhersagen, insbesondere eine natürliche Obergrenze für die Masse kompakter Kerne ($\sim 3 \times 10^6 M_\odot$), die Häufigkeit von Doppelkern-Systemen sowie fraktale Skalierungen in astrophysikalischen Plasmen. Sollten diese bestätigt werden, stünde ein Paradigmenwechsel in der Kosmologie bevor.

\begin{thebibliography}{99}

\bibitem{weber} Weber, W. (1846). Elektrodynamische Massbestimmungen. \textit{Annalen der Physik}, 143, 211-233.

\bibitem{bohm} Bohm, D. (1952). A Suggested Interpretation of the Quantum Theory in Terms of Hidden Variables. \textit{Physical Review}, 85, 166-193.

\bibitem{dstt} Czybor, M. (2025). Dynamische Schwere-Trägheits-Theorie (DSTT).

\bibitem{omega} Czybor, M. (2026). Die $\Omega$-Theorie der Gravitation.

\bibitem{fraktal} Czybor, M. (2025). Fraktale Maxwell-Gleichungen für Elektrodynamik und Gravitation.

\end{thebibliography}

\end{document}